\section{Methodology}

\subsection{Main Specification}

The main econometric exercise estimates the conditional firm-size wage premium in a log-wage model. The baseline specification is:
\begin{equation}\label{eq:main}
  \ln(w_{ist}) =
  \alpha
  + \sum_{g \neq g_0} \beta_g \mathbf{1}\{Size_{ist}=g\}
  + X_{ist}'\gamma
  + \lambda_{s \times t}
  + \varepsilon_{ist},    
\end{equation}
where \(w_{ist}\) is real hourly labor income for worker \(i\) in sector \(s\) in year \(t\), \(Size_{ist}\) is the harmonized firm-size category, and \(g_0\) is the reference group. We use solo workers as the reference group. This choice makes each coefficient directly interpretable as the conditional wage gap between workers in firm-size group \(g\) and workers at the bottom of the firm-size distribution. 

The vector \(X_{ist}\) includes a female-worker indicator, age, age squared, education dummies, and a formal-employment indicator. The fixed effect \(\lambda_{s \times t}\) absorbs sector-year conditions, allowing each broad sector to have its own wage level in each year. The model is estimated with GEIH expansion weights, and standard errors are clustered by sector.

The coefficients \(\beta_g\) are log-point differences relative to solo workers. For interpretation, we transform them into percentage premiums:
\begin{equation}
  Premium_g = 100 \times \left(\exp(\hat{\beta}_g)-1\right).    
\end{equation}

\subsection{Dynamic Premiums and Trend Test}

The same framework can be extended to study the evolution of the premium over time by interacting firm-size categories with year indicators:
\begin{equation}
  \ln(w_{ist}) =
  \alpha
  + \sum_{\tau}\sum_{g \neq g_0}
    \beta_{g\tau}\mathbf{1}\{Size_{ist}=g\}\mathbf{1}\{Year_{t}=\tau\}
  + X_{ist}'\gamma
  + \lambda_{s \times t}
  + \varepsilon_{ist}.    
\end{equation}

This dynamic specification produces a sequence of firm-size premiums by year and can be used to assess whether the premium has widened, narrowed, or remained stable. We also estimate a more parsimonious trend specification:
\begin{equation}
  \ln(w_{ist}) =
  \alpha
  + \sum_{g \neq g_0} \left(\beta_g
  +  \delta_g
    Trend_t\right)\mathbf{1}\{Size_{ist}=g\}
  + X_{ist}'\gamma
  + \lambda_{s \times t}
  + \varepsilon_{ist},
\end{equation}
where \(Trend_t\) is a linear year trend normalized to zero in 2008. The parameter \(\delta_g\) measures whether the wage premium for firm-size group \(g\) changed over time relative to solo workers. Positive values of \(\delta_g\) imply a steeper association between firm size and wages, while negative values imply a flatter one.


\subsection{Firm-Size Nonlinearities by Formality Status}

The descriptive evidence suggests that the relationship between firm size and
hourly income may differ between formal and informal workers. We therefore
estimate a specification that interacts all firm-size categories with formal
employment:

\begin{equation}\label{eq:formality-interaction}
  \ln(w_{ist}) =
  \alpha
  + \phi Formal_{ist}
  + \sum_{g \neq g_0} \left(\beta_g 
  + \theta_g
     Formal_{ist}\right)\mathbf{1}\{Size_{ist}=g\}
  + X_{ist}'\gamma
  + \lambda_{s \times t}
  + \varepsilon_{ist}.    
\end{equation}

The coefficients \(\beta_g\) measure the firm-size premium among informal workers, relative to informal solo workers. The corresponding premium among formal workers is \(\beta_g+\theta_g\), relative to formal solo workers. The coefficient \(\phi\) measures the formal-informal wage gap among solo workers, while \(\phi+\theta_g\) measures the formal-informal wage gap within firm-size group \(g\).

\subsection{Approximate Firm-Size Elasticities}

The main specifications treat firm size as a set of categories because this is
how firm size is reported in GEIH and because the largest category is top-coded.
This categorical approach is the most transparent way to estimate the Colombian
firm-size wage premium. However, much of the employer-size wage effect
literature reports a continuous elasticity of wages with respect to employer
size. To compare the Colombian evidence with those benchmarks, we also estimate
an approximate log-log specification:
\begin{equation}\label{eq:elasticity}
  \ln(w_{ist}) =
  \alpha
  + \eta \ln(\widetilde{Size}_{ist})
  + X_{ist}'\gamma
  + \lambda_{s \times t}
  + \varepsilon_{ist}.
\end{equation}

The variable \(\widetilde{Size}_{ist}\) assigns representative values to the
harmonized firm-size bins: 1 for solo workers, 2.5 for firms with 2--3 workers,
4.5 for 4--5 workers, 8 for 6--10 workers, 15 for 11--19 workers, 25 for
20--30 workers, 40.5 for 31--50 workers, 75.5 for 51--100 workers, and 150 for
the top-coded 101+ category. The coefficient \(\eta\) is therefore an
approximate elasticity: it measures the log-wage slope associated with assigned
firm size. Because the underlying variable is grouped and top-coded, these
elasticities are used only as benchmarking statistics. The categorical
specifications remain the main empirical estimates.
